\begin{abstractutfpr}
This work presents the development of \textit{Jump'n'Shoot}, a competitive three-dimensional multiplayer first-person shooter with arena combat and platforming elements. The objective was to investigate and integrate movement, combat, and network communication techniques by building a game with Godot Engine 4.5 and GDScript. The client handles input and presentation, while a dedicated server applies combat, resource, and scoring rules using positions reported by clients. The main technical contribution is the integration of these systems into a playable version featuring double jumps, jump pads, weapons, grenades, and collectible items. Teams compete in \textit{Team Deathmatch} and \textit{King of the Hill} on an arena designed for both modes. Synchronization uses remote procedure calls, interpolation of remote states, and historical state queries to validate shots. The source code is accompanied by documentation of the implementation decisions and limitations to support study and further development. In an exploratory test with the server and automated clients on the same machine, 300-second runs with eight, 16, and 32 participants recorded a mean round-trip time of approximately 19 ms and increasing processing time for the instrumented portion of the lobby update cycle as the number of participants increased.
\end{abstractutfpr}
